\section{A counting estimate}\label{sec:finite}

For a finitely supported sequence $u\in\mathbb R^\omega$ and
$\pi\in\Sym$, let
\[
 M(u)=\sup_{j<\omega}\left|S_j(u)\right| \quad \text{and}\quad
 M_\pi(u)=\sup_{j<\omega}\left|S_j^\pi(u)\right|.
\]
For $v\in\mathbb R^L$ and $\sigma\in\operatorname{Sym}(L)$, use
the same notation with the maximum over $j\leq L$.

When $v\in \{-p, p\}^L$ for some $p>0$, we think of $v$ as a walk starting at $0$.
We interpret $v_i=p$ as a step of length $p$ to the right and $v_i=-p$ as a step of length $p$ to the left.
Then $M(v)$ is the maximum distance from the origin reached by this walk.

We say that $v \in \{-p, p\}^L$ is \emph{balanced} if $S_L(v)=0$, that is, if the walk ends at the origin.
Clearly, balanced vectors exist only for even integers $L$.

\begin{lemma}\label{lem:finite}
For all positive integers $m$ and $q$,
there exist a positive integer $L$ and balanced vectors
$(v_k)_{k<m}$ in $\{-1/L,1/L\}^L$ satisfying, for every $\sigma\in\operatorname{Sym}(L)$,
\begin{equation*}
 \left|\{k<m:M_\sigma(v_k)>1/q\}\right|\leq 4q^4.
\end{equation*}
\end{lemma}

Recall that if $(e_k)_{k<m}$ is an orthonormal family in a real
inner product space, then Bessel's inequality gives
$\sum_{k<m}|\langle b,e_k\rangle|^2\leq\|b\|^2$ for every vector $b$
in that space.

\begin{proof}
List the elements of $\{-1,1\}^m$ as $(x_i)_{i<L}$, where $L=2^m$.
For $k<m$, put
\[
 v_k(i)=\frac{x_i(k)}L\qquad(i<L).
\]
Each $v_k$ is balanced, since
\begin{equation*}
      \sum_{i<L}v_k(i)=\frac1L\sum_{i<L}x_i(k)
      =\frac1L\sum_{x\in\{-1,1\}^m}x(k)=0.
\end{equation*}
Moreover, for every two distinct $k,\ell<m$,
\begin{equation}\label{eq:orthogonal}
 v_k\cdot v_\ell
 =\frac1{L^2}\sum_{x\in\{-1,1\}^m}x(k)x(\ell)
 =\frac1{L^2}\left(\sum_{x(k)=x(\ell)}1-\sum_{x(k)\neq x(\ell)}1\right)
 =0.
\end{equation}
Also, for every $k<m$,
\begin{equation}\label{eq:norm}
 \|v_k\|^2=\sum_{i<L}v_k(i)^2
 =\sum_{i<L}\frac1{L^2}=\frac1L.
\end{equation}

Fix $\sigma\in\operatorname{Sym}(L)$.
The vectors $e_k=\sqrt L\,(v_k\circ\sigma)$ are orthonormal by \eqref{eq:orthogonal} and \eqref{eq:norm}.
For $j\leq L$, define $b_j\in\mathbb R^L$ by
\[
 b_j(i)=
 \begin{cases}
  1,&i<j,\\
  0,&j\leq i<L.
 \end{cases}
\]
Then, for each $j\leq L$,
\begin{equation}\label{eq:prefix-square}
 \sum_{k<m}|S_j^\sigma(v_k)|^2
 =\frac1L\sum_{k<m}|\langle b_j,e_k\rangle|^2
 \leq\frac{\|b_j\|^2}{L}
 =\frac jL\leq1.
\end{equation}

For $k<m$ and $j\leq L$, telescoping gives
\[
\begin{aligned}
 |S_j^\sigma(v_k)|^2
 &=\sum_{i<j}\left((S_{i+1}^\sigma(v_k))^2
                         -(S_i^\sigma(v_k))^2\right)\\
 &=\sum_{i<j}\left(S_{i+1}^\sigma(v_k)-S_i^\sigma(v_k)\right)
             \left(S_{i+1}^\sigma(v_k)+S_i^\sigma(v_k)\right)\\
 &=\sum_{i<j}v_k(\sigma(i))
             \left(S_{i+1}^\sigma(v_k)+S_i^\sigma(v_k)\right)\\
 &\leq\frac1L\sum_{i<L}
             \left(|S_{i+1}^\sigma(v_k)|+|S_i^\sigma(v_k)|\right)
 =\frac2L\sum_{i<L}|S_i^\sigma(v_k)|,
\end{aligned}
\]
where the last equality uses $S_0^\sigma(v_k)=S_L^\sigma(v_k)=0$.
Taking the maximum over $j$ and applying the Cauchy--Schwarz inequality, we get
\[
 M_\sigma(v_k)^4
 \leq\left(\frac2L\sum_{i<L}1\cdot|S_i^\sigma(v_k)|\right)^2
 \leq\frac4L\sum_{i<L}|S_i^\sigma(v_k)|^2.
\]
Summing over $k$ and using \eqref{eq:prefix-square},
\[
 \sum_{k<m}M_\sigma(v_k)^4
 \leq\frac4L\sum_{i<L}\sum_{k<m}|S_i^\sigma(v_k)|^2
 \leq\frac4L\sum_{i<L}1=4.
\]
Finally, let $B=\{k<m:M_\sigma(v_k)>1/q\}$.
The required bound follows from
\[
 \frac{|B|}{q^4}=\sum_{k\in B}\frac1{q^4}
 \leq\sum_{k\in B}M_\sigma(v_k)^4
 \leq\sum_{k<m}M_\sigma(v_k)^4\leq4.\qedhere
\]
\end{proof}
